\pg{429} \emph{operações} dos agentes \emph{Zrraciomaes}, --- os \emph{moraes}, na existencia da lei \emph{moral} e no \emph{consenso} de todos os povos \footnote{Na classificação dos argumentos, que demonstram a existencia de \textsc{Deus}, em \emph{methaphysicos}, \emph{physicos}, \emph{e moraes}, conforme se baséam na \emph{essencia} das coisas, ou na \emph{ordem} existente no mundo e nas \emph{operações} dos agentes \emph{irracionaes}, ou na lei \emph{moral} e no \emph{consenso} dos povos, seguimos a doutrina dos mais sensatos auctores.\par Alguns dividem esses argumentos em \emph{cosmologicos}, \emph{psychologicos} e \emph{moraes}; conforme os factos, em que se fundam, são \emph{cosmologicos} (baseados nas leis \emph{physicas}), \emph{psychologicos} (baseados nas leis \emph{intellectuaes}), ou \emph{moraes} (baseados nas leis \emph{moraes}). --- Esta divisão tambem é acceitavel.\par \emph{Kant}, dizendo que o espirito humano se eleva a \textsc{DEus} por meio de um d'estes tres principios: de \emph{identidade}, de \emph{causalidade} ou de \emph{finalidade}, reduz as provas a tres classes: \emph{ontologicas}, \emph{cosmologicas} e \emph{physico-teleologicas}, conforme se baseam no principio de \emph{identidade}, de \emph{causalidade} ou de \emph{finalidade}. --- A. doutrina de \emph{Kant} é inacceitavel. O principio, em que se baseam todas as provas em favor da existencia de \textsc{Deus}, é um só: o \emph{principio de causalidade}, e a distineção das provas deriva da natureza dos \emph{factos}, aos quaes o dito principio se applica. O principio de \emph{identidade} não prova coisa alguma, como o proprio \emph{Kant} confessa; o principio de \emph{finalidade}, longe de ser a base do argumento \emph{fteleologíco}, é apenas um corollario d'esse argumento.\par A. existencia de \textsc{DEus} é uma verdade tão necessaria e fundamental, que recorre em todos os grandes problemas da philosophia. Nem admira. Nào encontrando nos agentes finitos a causa \emph{sufficiente} dos seres, a razão deve elevar-se a uma causa \emph{suprema} e \emph{infinita}. \emph{A origem} do \emph{mundo} (\emph{Cosm}., \emph{ger}.\pgn{430}Cap. I, art. V), á \emph{finalidade} nas operações dos seres irracionaes (\emph{Cosm}. \emph{ger}., Cap. II, art. III a ordem \emph{physica} com as suas \emph{leis} (\emph{Cosm}. \emph{ger}., Cap. II, art. IV e V), a \emph{origem} e o \emph{desenvolvimento da vida} (\emph{Cosm}. \emph{esp}., Cap. II, art. III), á \emph{origem} da \emph{alma} e do \emph{corpo} humano (\emph{Anthr}., Sec. IV, Cap. D, tudo o que nos cerca, com os seus caracteres de \emph{contingencia} e de \emph{mutabilidade}, proclama a existencia de uma Causa infinitamente sabia, poderosa e boa --- a existencia de \textsc{DEus}.\par Escutámos esta voz maravilhosa da creação e podemos dizer, com immensa satisfação da nossa alma, que procurámos sempre pôr em relevo a racionabilidade e a necessidade da existencia de um Ente supremo, e mostrar a fatuidade, os erros e as contradieções dos adversarios.\par Mas isso não obsta a que façamos aqui, como no logar proprio, uma synthese do que dissemos em todo este trabalho, indieando os principaes argumentos em defeza da \emph{existencia} de Deus. --- Cf. \emph{Sum}. \emph{Th}., p. I, q. 2, a. 8; --- \emph{c}. \emph{Gentes}, l. I, c. 18.}.

. 12. Argumentos \textbf{metaphysicos}. --- Os argumentos \emph{metaphysicos}, em favor da \emph{existencia} de \textsc{Deus}, baséam-se em que, --- a) se no mundo existe o movimento ou mudança, deve existir um motor \emph{emmovel}, --- \emph{b}) se existem as causas \emph{segundas}, deve existir a causa \emph{primeira}, --- c) se existem seres \emph{contingentes}, deve existir um ser \emph{necessario}, --- \emph{d}) se existem varios e diversos graus de perfeição, deve existir um ser infinitamente perfeito, --- e) se existem verdades \emph{absolutas} e \emph{necessarias}, deve existir uma intelligencia \emph{absoluta} e \emph{necessaria}.

